
This work presented a systematic study of cross-family probe transfer for hallucination detection. Probes trained on Llama-3-8B transfer to Mistral-7B and Qwen-2-7B hidden states with a mean AUROC gap of only 0.0133.

The main contributions are:

\begin{enumerate}
\item \textbf{First systematic cross-family validation.} SEP transfer was tested across three distinct LLM families (Meta, Mistral AI, Alibaba), demonstrating that all six cross-family pairs achieve transfer gaps below 0.034.
\end{enumerate}

\begin{enumerate}
\item \textbf{Affine alignment for dimension mismatch.} Least-squares affine mapping enables transfer between models with different hidden dimensions, with gaps as small as 0.003.
\end{enumerate}

\begin{enumerate}
\item \textbf{Evidence for convergent uncertainty encoding.} The transfer matrix provides empirical support for the hypothesis that uncertainty encoding shares geometric properties across transformer architectures.
\end{enumerate}

\textbf{Future Directions.} Extending this methodology to 70B+ models would determine whether the architecture-invariant encoding persists at scale. Systematic evaluation across additional benchmarks (HaluEval, NQ-Open) would establish the generality of the findings. Finally, full evaluation with true semantic entropy labels (rather than proof-of-concept random labels) is needed to confirm absolute performance levels.
